\chapter{Inner Products and Projections}\label{linear:ch:inner}

\begin{center}\itshape
Of all the principles that can be proposed for this purpose, I think there is none more general, more exact, or easier
to apply than the one we have used in this work: it consists of making the sum of the squares of the errors a minimum.
\upshape(Adrien-Marie Legendre, \emph{Nouvelles m\'ethodes pour la d\'etermination des orbites des com\`etes}, 1805)
\end{center}

On the first night of the nineteenth century, 1 January 1801, Giuseppe Piazzi in Palermo noticed a faint star that
was not on his charts. The next night it had moved. He followed it for six weeks, through about three degrees of sky,
until it vanished into the glare of the Sun. Piazzi had found Ceres, the first asteroid, and almost immediately lost
it: from six weeks of noisy positions, nobody could say where it would reappear in the autumn. The astronomers of
Europe published predictions, and they disagreed.

A twenty-four-year-old in Brunswick, Carl Friedrich Gauss, set aside his number theory and attacked the problem. His
prediction differed from everyone else's, by several degrees. On 31 December 1801, almost exactly a year after Piazzi's
discovery, Franz von Zach pointed his telescope where Gauss said, and Ceres was there. Gauss became famous overnight.
He did not explain his method until 1809, and by then Legendre had published, in 1805, the principle quoted above: when
a system has more equations than unknowns and no exact solution, choose the unknowns that make the sum of the squared
errors smallest. Gauss wrote that he had used this principle since 1795. Legendre was not amused, and the priority
quarrel that followed is one of the classic bad tempers of the history of mathematics.

The geometry behind least squares is the subject of this chapter. Once a vector space has an \emph{inner product}, it
has lengths and angles, and the closest point of a subspace to a given vector is found by dropping a perpendicular. That
single picture gives the Cauchy--Schwarz inequality, Gram--Schmidt and the $QR$ factorisation, least squares, Gram
determinants as squared volumes, and, applied to the space of polynomials, the orthogonal polynomials and Gauss's most
beautiful numerical discovery, a quadrature rule that is exact twice as often as it has any right to be.

\begin{think}
The three points $(0,0)$, $(1,1)$, $(2,1)$ are not on a line. Which line $y=a+bx$ is the ``closest'' to them, in the
sense that the sum of the squares of the vertical errors is smallest? Before calculating, guess whether the line passes
above or below the middle point, and whether the errors have any reason to add up to zero. Example~\ref{linear:exa:line} answers.
\end{think}

\section{Inner products and the Cauchy--Schwarz inequality}

Throughout, $\mathbb K$ is $\R$ or $\C$, and $\bar z$ is the complex conjugate (so $\bar z=z$ when $\mathbb K=\R$).

\begin{definition}\label{linear:def:inner}
An \emph{inner product} on a $\mathbb K$-vector space $V$ is a map $V\times V\to\mathbb K$, $(x,y)\mapsto\langle x,y\rangle$,
that is linear in $x$, satisfies $\langle y,x\rangle=\overline{\langle x,y\rangle}$, and is \emph{positive definite}:
$\langle x,x\rangle>0$ for $x\ne0$. The \emph{norm} is $\|x\|=\sqrt{\langle x,x\rangle}$; vectors with $\langle x,y\rangle=0$
are \emph{orthogonal}, written $x\perp y$.
\end{definition}

It follows that $\langle x,\lambda y+\mu z\rangle=\bar\lambda\langle x,y\rangle+\bar\mu\langle x,z\rangle$: the inner
product is conjugate-linear in the second variable. The standard examples are $\langle x,y\rangle=\sum_ix_i\bar y_i=y^*x$
on $\mathbb K^n$ (where $y^*=\bar y^{\mathsf T}$), and
\[
\langle f,g\rangle=\int_a^bf(t)\overline{g(t)}\,w(t)\dd t
\]
on continuous functions on $[a,b]$, or on polynomials, for a positive continuous \emph{weight} $w$. Positivity holds
because a continuous function with $\int_a^b|f|^2w=0$ vanishes identically.

The expansion $\|x+y\|^2=\|x\|^2+2\operatorname{Re}\langle x,y\rangle+\|y\|^2$ is used constantly; for orthogonal $x,y$
it is \emph{Pythagoras}: $\|x+y\|^2=\|x\|^2+\|y\|^2$.

\begin{theorem}[Cauchy--Schwarz]\label{linear:thm:cs}
For all $x,y\in V$, $|\langle x,y\rangle|\le\|x\|\,\|y\|$, with equality if and only if $x$ and $y$ are linearly dependent.
\end{theorem}
\begin{proof}
If $y=0$ both sides vanish. Otherwise let $p=\frac{\langle x,y\rangle}{\|y\|^2}\,y$, the ``shadow'' of $x$ on the line
$\mathbb Ky$. Then
\[
\langle x-p,y\rangle=\langle x,y\rangle-\frac{\langle x,y\rangle}{\|y\|^2}\langle y,y\rangle=0,
\]
so $x-p\perp p$ and Pythagoras gives $\|x\|^2=\|p\|^2+\|x-p\|^2\ge\|p\|^2=|\langle x,y\rangle|^2/\|y\|^2$. Equality holds
exactly when $x=p$, that is when $x\in\mathbb Ky$.
\end{proof}

The proof is the whole chapter in miniature: drop a perpendicular, then apply Pythagoras. Two consequences follow.

\begin{corollary}\label{linear:cor:triangle}
$\|x+y\|\le\|x\|+\|y\|$, so $d(x,y)=\|x-y\|$ is a distance. In a real inner product space, the \emph{angle}
$\theta\in[0,\pi]$ between non-zero vectors, defined by $\cos\theta=\frac{\langle x,y\rangle}{\|x\|\|y\|}$, makes sense.
\end{corollary}
\begin{proof}
$\|x+y\|^2=\|x\|^2+2\operatorname{Re}\langle x,y\rangle+\|y\|^2\le\|x\|^2+2\|x\|\|y\|+\|y\|^2$. The quotient defining
$\cos\theta$ lies in $[-1,1]$ by Theorem~\ref{linear:thm:cs}.
\end{proof}

Applied to the examples, Cauchy--Schwarz says $\bigl|\sum x_iy_i\bigr|^2\le\sum x_i^2\sum y_i^2$ and
$\bigl|\int fg\bigr|^2\le\int f^2\int g^2$: two inequalities of analysis that are one inequality of geometry.

\begin{remark}
Not every norm comes from an inner product: a norm does if and only if it satisfies the \emph{parallelogram law}
$\|x+y\|^2+\|x-y\|^2=2\|x\|^2+2\|y\|^2$ (Jordan and von Neumann, 1935). One direction is a two-line expansion; the other,
harder one is Exercise~\ref{linear:ex:jvn}. The norm $|x_1|+|x_2|$ on $\R^2$ fails the law at $x=(1,0)$, $y=(0,1)$.
\end{remark}

\section{Orthonormal bases and Gram--Schmidt}

A family $(e_1,\dots,e_k)$ is \emph{orthonormal} if $\langle e_i,e_j\rangle=\delta_{ij}$.

\begin{proposition}\label{linear:prop:orthfamily}
An orthonormal family is linearly independent, and for $x=\sum_i x_ie_i$ in its span, $x_i=\langle x,e_i\rangle$ and
$\|x\|^2=\sum_i|x_i|^2$.
\end{proposition}
\begin{proof}
Taking the inner product of $x=\sum_jx_je_j$ with $e_i$ gives $\langle x,e_i\rangle=x_i$; in particular $x=0$ forces
every $x_i=0$. Then $\|x\|^2=\langle\sum_ix_ie_i,\sum_jx_je_j\rangle=\sum_{i,j}x_i\bar x_j\delta_{ij}=\sum_i|x_i|^2$.
\end{proof}

So in an orthonormal basis, coordinates are computed by inner products, with no linear system to solve: this is what
makes orthonormal bases so convenient, and the next theorem says they always exist.

\begin{theorem}[Gram--Schmidt]\label{linear:thm:gramschmidt}
Let $v_1,\dots,v_k$ be linearly independent in an inner product space. There is a unique orthonormal family
$e_1,\dots,e_k$ such that, for each $j$, $\operatorname{span}(e_1,\dots,e_j)=\operatorname{span}(v_1,\dots,v_j)$ and
$\langle v_j,e_j\rangle>0$. Explicitly,
\[
u_j=v_j-\sum_{i<j}\langle v_j,e_i\rangle e_i,\qquad e_j=\frac{u_j}{\|u_j\|}.
\]
\end{theorem}
\begin{proof}
By induction on $j$. Suppose $e_1,\dots,e_{j-1}$ are built, orthonormal, with span $W_{j-1}=\operatorname{span}(v_1,\dots,
v_{j-1})$. The vector $u_j$ satisfies $\langle u_j,e_i\rangle=\langle v_j,e_i\rangle-\langle v_j,e_i\rangle=0$ for $i<j$,
and $u_j\ne0$ since $v_j\notin W_{j-1}$. So $e_j=u_j/\|u_j\|$ extends the orthonormal family, $\operatorname{span}(e_1,
\dots,e_j)=W_{j-1}+\mathbb Kv_j$, and $\langle v_j,e_j\rangle=\langle u_j,e_j\rangle=\|u_j\|>0$ (the difference $v_j-u_j$
lies in $W_{j-1}\perp e_j$).

Uniqueness: any $e_j'$ with the stated properties lies in $W_j$ and is orthogonal to $W_{j-1}$, a hyperplane of $W_j$; so
$e_j'$ is a unit multiple $ce_j$, and $\langle v_j,e_j'\rangle=\bar c\|u_j\|>0$ forces $c=1$.
\end{proof}

In particular every finite-dimensional inner product space has an orthonormal basis, and every orthonormal family extends
to one (apply Gram--Schmidt to the family followed by any basis, discarding vectors that become dependent). Written with
matrices, Gram--Schmidt is a factorisation.

\begin{corollary}[$QR$ factorisation]\label{linear:cor:qr}
Let $A\in M_{m,n}(\mathbb K)$ have linearly independent columns. Then $A=QR$, where $Q\in M_{m,n}(\mathbb K)$ has orthonormal
columns ($Q^*Q=I_n$) and $R\in M_n(\mathbb K)$ is upper triangular with positive diagonal entries. This factorisation is
unique.
\end{corollary}
\begin{proof}
Apply Theorem~\ref{linear:thm:gramschmidt} to the columns $a_j$ of $A$. Since $a_j\in\operatorname{span}(e_1,\dots,e_j)$,
Proposition~\ref{linear:prop:orthfamily} gives $a_j=\sum_{i\le j}r_{ij}e_i$ with $r_{ij}=\langle a_j,e_i\rangle$ and
$r_{jj}=\|u_j\|>0$; this is $A=QR$ with $Q=(e_1|\cdots|e_n)$.

If also $A=Q'R'$, the columns of $Q$ and of $Q'$ are orthonormal bases of the same space, the column space of $A$.
Then $U=Q'^*Q$ satisfies $Q=Q'U$ (expand each column of $Q$ in the basis given by $Q'$), so $U$ is the matrix of a change
between orthonormal bases and $U^*U=Q^*Q'Q'^*Q=Q^*Q=I$. From $Q'R'=QR=Q'UR$ we get $R'=UR$, so $U=R'R^{-1}$ is upper
triangular with positive diagonal. Its inverse $U^*$ is then both upper triangular (as an inverse of an upper triangular
matrix) and lower triangular (as a conjugate transpose), so $U$ is diagonal, with entries of modulus $1$ that are
positive: $U=I$.
\end{proof}

\section{Orthogonal projections}

For a subset $S$ of $V$, the \emph{orthogonal} $S^\perp=\{x:\ x\perp s\text{ for all }s\in S\}$ is a subspace.

\begin{theorem}[orthogonal projection]\label{linear:thm:projection}
Let $W$ be a finite-dimensional subspace of an inner product space $V$ (which may have infinite dimension), with
orthonormal basis $e_1,\dots,e_k$, and let
\[
P_Wx=\sum_{i=1}^k\langle x,e_i\rangle e_i.
\]
Then $x-P_Wx\in W^\perp$, so $V=W\oplus W^\perp$. The point $P_Wx$ is the unique closest point of $W$ to $x$: for
$w\in W$, $\|x-w\|^2=\|x-P_Wx\|^2+\|P_Wx-w\|^2$. The map $P_W$ is linear, $P_W^2=P_W$, and
$\langle P_Wx,y\rangle=\langle x,P_Wy\rangle$.
\end{theorem}
\begin{proof}
For each $j$, $\langle x-P_Wx,e_j\rangle=\langle x,e_j\rangle-\langle x,e_j\rangle=0$, so $x-P_Wx\perp W$. Hence
$x=P_Wx+(x-P_Wx)$ with the first term in $W$ and the second in $W^\perp$; and $W\cap W^\perp=0$ because a vector
orthogonal to itself is zero. For $w\in W$, the vector $P_Wx-w$ lies in $W$, so it is orthogonal to $x-P_Wx$, and
Pythagoras applied to $x-w=(x-P_Wx)+(P_Wx-w)$ gives the displayed identity: $\|x-w\|$ is smallest exactly when $w=P_Wx$.
Linearity is clear, $P_W$ fixes $W$ so $P_W^2=P_W$, and $\langle P_Wx,y\rangle=\sum_i\langle x,e_i\rangle\overline{\langle
y,e_i\rangle}$ is symmetric in the required sense.
\end{proof}

\begin{figure}[ht]
\centering
\fig{linear-ch04-projection}
\caption{The orthogonal projection. Among all points $w$ of the plane $W$, the foot $P_Wx$ of the perpendicular is the
closest to $x$, because $x-w$ is the hypotenuse of a right triangle whose other side is $x-P_Wx$.}
\label{linear:fig:projection}
\end{figure}

Since $\|x\|^2=\|P_Wx\|^2+\|x-P_Wx\|^2$, we obtain \emph{Bessel's inequality}: for every orthonormal family $(e_i)$,
\begin{equation}\label{linear:eq:bessel}
\sum_i|\langle x,e_i\rangle|^2\le\|x\|^2.
\end{equation}
It holds for infinite orthonormal families too, since every finite partial sum is bounded. Equality for every $x$ is
Parseval's identity, the subject of Hilbert spaces (Tome~7, Chapter~2); Exercise~\ref{linear:ex:basel} squeezes a famous
inequality out of Bessel alone.

In finite dimension, $W^{\perp\perp}=W$ and $\dim W^\perp=\dim V-\dim W$, which is the inner-product version of the
annihilator of Chapter~\ref{linear:ch:duality} (Theorem~\ref{linear:thm:annihilator}): the inner product identifies $V$ with its dual,
as we shall see in Section~\ref{linear:sec:adjoint}.

\section{Least squares}

Gauss's problem in modern clothes: $A\in M_{m,n}(\mathbb K)$ with $m>n$, $b\in\mathbb K^m$, and the system $Ax=b$ has no
solution. We look for $x$ minimising $\|Ax-b\|$.

\begin{theorem}[least squares]\label{linear:thm:lsq}
The vector $x$ minimises $\|Ax-b\|$ if and only if it satisfies the \emph{normal equations}
\[
A^*Ax=A^*b.
\]
Minimisers always exist; the minimiser is unique if and only if the columns of $A$ are independent, and then
$x=(A^*A)^{-1}A^*b=R^{-1}Q^*b$, where $A=QR$.
\end{theorem}
\begin{proof}
As $x$ varies, $Ax$ runs over $W=\operatorname{Im}A$, and by Theorem~\ref{linear:thm:projection} the closest point of $W$ to
$b$ is $P_Wb$. So $x$ is a minimiser if and only if $Ax=P_Wb$, if and only if $b-Ax\perp W$, that is
$\langle b-Ax,Ay\rangle=y^*A^*(b-Ax)=0$ for every $y$, which is $A^*(b-Ax)=0$. Minimisers exist because $P_Wb\in W$.
They are unique if and only if $\ker A=0$. In that case $A^*A$ is invertible, because $A^*Ax=0$ implies
$\|Ax\|^2=x^*A^*Ax=0$. With $A=QR$, the normal equations read $R^*Q^*QRx=R^*Q^*b$, that is $R^*Rx=R^*Q^*b$, and $R^*$ is
invertible.
\end{proof}

\begin{example}\label{linear:exa:line}
For the three points of the opening question, the unknowns are $(a,b)$ and
\[
A=\begin{pmatrix}1&0\\1&1\\1&2\end{pmatrix},\quad y=\begin{pmatrix}0\\1\\1\end{pmatrix},\quad
A^{\mathsf T}A=\begin{pmatrix}3&3\\3&5\end{pmatrix},\quad A^{\mathsf T}y=\begin{pmatrix}2\\3\end{pmatrix}.
\]
The normal equations $3a+3b=2$, $3a+5b=3$ give $b=\frac12$, $a=\frac16$. The errors $y_i-(a+bx_i)$ are $-\frac16,\frac13,
-\frac16$: the line passes \emph{below} the middle point, and the errors add up to zero. That is no accident: the first
normal equation says that the error vector is orthogonal to the first column of $A$, the vector $(1,1,1)$. Whenever
the model contains a constant term, least squares errors sum to zero.
\end{example}

\begin{figure}[ht]
\centering
\fig{linear-ch04-lsq}
\caption{The least squares line of Example~\ref{linear:exa:line}. The sum of the squares of the red segments is $\frac16$,
smaller than for any other line.}
\label{linear:fig:lsq}
\end{figure}

Why squares, and not absolute values or fourth powers? Legendre's answer was convenience: squares give linear equations.
Gauss's answer, given in 1821, is deeper and needs only the vocabulary of expectation and covariance (Tome~3). Suppose
that $b=Ax_0+\varepsilon$, where $x_0$ is the unknown truth and the noise $\varepsilon$ is a random vector with
$\mathbb E\varepsilon=0$ and covariance $\mathbb E\varepsilon\varepsilon^{\mathsf T}=\sigma^2I$ (independent errors of
equal size suffice), all real. A \emph{linear unbiased estimator} is $\hat x=Cb$ with $\mathbb E\hat x=x_0$ for every
$x_0$, that is $CA=I$; its covariance matrix is $\mathbb E(\hat x-x_0)(\hat x-x_0)^{\mathsf T}=\sigma^2CC^{\mathsf T}$.

\begin{theorem}[Gauss--Markov]\label{linear:thm:gaussmarkov}
Let $A$ have independent columns and $C_0=(A^{\mathsf T}A)^{-1}A^{\mathsf T}$, the least squares estimator. For every
$C$ with $CA=I$, the matrix $CC^{\mathsf T}-C_0C_0^{\mathsf T}$ is positive semidefinite. In particular, for every
vector $u$, the variance of $u^{\mathsf T}Cb$ is at least that of $u^{\mathsf T}C_0b$: least squares is the best linear
unbiased estimator.
\end{theorem}
\begin{proof}
Let $D=C-C_0$. Then $DA=CA-C_0A=I-I=0$, hence $DC_0^{\mathsf T}=DA(A^{\mathsf T}A)^{-1}=0$ and also $C_0D^{\mathsf
T}=(DC_0^{\mathsf T})^{\mathsf T}=0$. Expanding,
\[
CC^{\mathsf T}=(C_0+D)(C_0+D)^{\mathsf T}=C_0C_0^{\mathsf T}+DD^{\mathsf T},
\]
and $u^{\mathsf T}DD^{\mathsf T}u=\|D^{\mathsf T}u\|^2\ge0$. The variance of $u^{\mathsf T}Cb$ is
$\sigma^2u^{\mathsf T}CC^{\mathsf T}u$.
\end{proof}

The proof is once more a right angle: the competitor $C$ is the least squares estimator plus a perpendicular error
$D$, and Pythagoras says the error can only add variance.

\section{Gram determinants and volumes}

The \emph{Gram matrix} of vectors $v_1,\dots,v_k$ is $G=(\langle v_j,v_i\rangle)_{i,j}$; if the $v_j$ are the columns of
a matrix $A\in M_{m,k}(\mathbb K)$ with the standard inner product, then $G=A^*A$. Its determinant is the \emph{Gram
determinant} $\operatorname{Gram}(v_1,\dots,v_k)$.

\begin{theorem}[Gram determinants]\label{linear:thm:gram}
Let $u_j=v_j-P_{W_{j-1}}v_j$, where $W_{j-1}=\operatorname{span}(v_1,\dots,v_{j-1})$: the part of $v_j$ perpendicular to
the earlier vectors, as in Gram--Schmidt.
\begin{enumerate}[label=(\alph*)]
\item $\operatorname{Gram}(v_1,\dots,v_k)=\|u_1\|^2\|u_2\|^2\cdots\|u_k\|^2$. In particular the Gram determinant is
$\ge0$, and it vanishes if and only if the $v_j$ are dependent.
\item If $v_1,\dots,v_k$ are independent and $W=\operatorname{span}(v_1,\dots,v_k)$, the distance from $x$ to $W$ is given
by
\[
d(x,W)^2=\frac{\operatorname{Gram}(v_1,\dots,v_k,x)}{\operatorname{Gram}(v_1,\dots,v_k)}.
\]
\item For $v_1,\dots,v_k\in\R^n$, the columns of a matrix $A$,
\[
\operatorname{Gram}(v_1,\dots,v_k)=\sum_{|I|=k}\det(A_I)^2,
\]
the sum running over the $k\times k$ minors $A_I$ made of the rows in $I$.
\end{enumerate}
\end{theorem}
\begin{proof}
(a) If the $v_j$ are dependent, some $c\ne0$ has $Ac=0$ (working in coordinates in an orthonormal basis of their span),
so $Gc=A^*Ac=0$ and $\det G=0$; and some $u_j$ is zero. If they are independent, write $A=QR$ (Corollary~\ref{linear:cor:qr});
then $G=A^*A=R^*Q^*QR=R^*R$ and $\det G=|\det R|^2=\prod_jr_{jj}^2=\prod_j\|u_j\|^2$, since $r_{jj}=\|u_j\|$.
(b) Apply (a) to $(v_1,\dots,v_k,x)$: the last factor is $\|x-P_Wx\|^2=d(x,W)^2$, and the others are those of (a).
(c) is the Pythagoras theorem for $k$-volumes, Corollary~\ref{linear:cor:pythagoras}, since $G=A^{\mathsf T}A$.
\end{proof}

Part (a) says that $\sqrt{\operatorname{Gram}}$ is ``base times height'' iterated: the $k$-dimensional volume of the
parallelepiped spanned by the $v_j$, as in Chapter~\ref{linear:ch:volume}. For $k=2$ it is Cauchy--Schwarz again:
$\|x\|^2\|y\|^2-|\langle x,y\rangle|^2\ge0$. Part (b) turns any best-approximation problem into a ratio of determinants,
and the next section has a family of determinants we can compute exactly.

\section{The Cauchy determinant and the Hilbert matrix}

\begin{theorem}[Cauchy, 1841]\label{linear:thm:cauchydet}
For numbers $x_1,\dots,x_n$ and $y_1,\dots,y_n$ with all $x_i+y_j\ne0$,
\[
\det\Bigl(\frac1{x_i+y_j}\Bigr)_{i,j=1}^n=\frac{\prod_{i<j}(x_j-x_i)(y_j-y_i)}{\prod_{i,j}(x_i+y_j)}.
\]
\end{theorem}
\begin{proof}
By induction on $n$; the case $n=1$ is clear. Call the matrix $C$. Subtract the last row from each other row: for $i<n$,
\[
\frac1{x_i+y_j}-\frac1{x_n+y_j}=\frac{x_n-x_i}{(x_i+y_j)(x_n+y_j)}.
\]
Factor $x_n-x_i$ out of row $i$ (for $i<n$) and $\frac1{x_n+y_j}$ out of column $j$ (for all $j$). The rows $i<n$ become
$\frac1{x_i+y_j}$ again, and the last row becomes $(1,\dots,1)$. Now subtract the last column from each other column:
for $i<n$ and $j<n$,
\[
\frac1{x_i+y_j}-\frac1{x_i+y_n}=\frac{y_n-y_j}{(x_i+y_j)(x_i+y_n)},
\]
and the last row becomes $(0,\dots,0,1)$. Factor $y_n-y_j$ out of column $j<n$ and $\frac1{x_i+y_n}$ out of row $i<n$,
and expand along the last row. What remains is the Cauchy matrix of $x_1,\dots,x_{n-1}$ and $y_1,\dots,y_{n-1}$, and
\[
\det C=\frac{\prod_{i<n}(x_n-x_i)\prod_{j<n}(y_n-y_j)}{\prod_{j\le n}(x_n+y_j)\prod_{i<n}(x_i+y_n)}\cdot\det C_{n-1}.
\]
The fraction contains exactly the factors of the formula that involve the index $n$, so induction concludes.
\end{proof}

In 1894 Hilbert asked how small $\int_0^1P(t)^2\dd t$ can be for a non-zero polynomial $P$ with integer coefficients. Its
answer goes through the Gram matrix of $1,t,\dots,t^{n-1}$ in $L^2[0,1]$, the \emph{Hilbert matrix}
\[
H_n=\Bigl(\int_0^1t^{i+j}\dd t\Bigr)_{i,j=0}^{n-1}=\Bigl(\frac1{i+j+1}\Bigr)_{i,j=0}^{n-1}
=\begin{pmatrix}1&\frac12&\frac13&\cdots\\[2pt]\frac12&\frac13&\frac14&\cdots\\[2pt]\frac13&\frac14&\frac15&\cdots\\[-2pt]
\vdots&\vdots&\vdots&\ddots\end{pmatrix}.
\]
Being a Gram matrix of independent vectors, $H_n$ is positive definite, a fact that would be painful to check directly.

\begin{corollary}\label{linear:cor:hilbert}
With $c_n=\prod_{k=1}^{n-1}k!$, we have $\det H_n=\dfrac{c_n^4}{c_{2n}}$, and in $L^2[0,1]$ the distance from $t^n$ to the
polynomials of degree less than $n$ is
\[
d\bigl(t^n,\R[t]_{<n}\bigr)=\frac{(n!)^2}{(2n)!\sqrt{2n+1}}=\frac1{\binom{2n}n\sqrt{2n+1}}.
\]
\end{corollary}
\begin{proof}
$H_n$ is the Cauchy matrix of Theorem~\ref{linear:thm:cauchydet} with $x_i=i$ and $y_j=j+1$ for $i,j=0,\dots,n-1$.
Passing from $H_n$ to $H_{n+1}$ adds $x_n=n$ and $y_n=n+1$, and multiplies the determinant by the factors of Cauchy's
formula that involve the index $n$:
\[
\frac{\det H_{n+1}}{\det H_n}=\frac{\prod_{i<n}(n-i)^2}{\prod_{j\le n}(n+j+1)\prod_{i<n}(i+n+1)}
=\frac{(n!)^2}{\frac{(2n+1)!}{n!}\cdot\frac{(2n)!}{n!}}=\frac{(n!)^4}{(2n)!\,(2n+1)!}.
\]
Since $c_{n+1}=c_n\,n!$ and $c_{2n+2}=c_{2n}\,(2n)!\,(2n+1)!$, the quotient $\frac{c_{n+1}^4}{c_{2n+2}}\big/\frac{c_n^4}{c_{2n}}$
is the same number, and $\det H_1=1=c_1^4/c_2$; the formula for $\det H_n$ follows by induction. By Theorem~\ref{linear:thm:gram}(b), the squared distance is the ratio
$\det H_{n+1}/\det H_n=(n!)^4/\bigl((2n)!(2n+1)!\bigr)$, which is the square of the stated value.
\end{proof}

So $t$ is at distance $\frac1{2\sqrt3}$ from the constants (the closest constant is $\frac12$), $t^2$ at distance
$\frac1{6\sqrt5}$ from the lines, and $t^n$ at distance about $\frac{\sqrt{\pi n}}{4^n\sqrt{2n+1}}$ from lower degrees: the
monomials become almost dependent very fast. Correspondingly $\det H_n$ is tiny, about $4^{-n^2}$, and $H_n$ is a
classic torture test for numerical software . Meanwhile $1/\det H_n$ is an integer
(Exercise~\ref{linear:ex:hilbertint}).

\begin{remark}
Hilbert's application: for an interval of length less than $4$, there are non-zero integer polynomials with
$\int P^2$ as small as we like, because $\det H_n$ for that interval tends to $0$ fast enough for Minkowski's lattice point
theorem to apply (Tome~12, Chapter~3, proves Minkowski's theorem). For the general weighted version, and the open problem
of the exact threshold for the uniform norm (the \emph{integer Chebyshev constant}), see the article of Borwein and
Erd\'elyi in the sources.
\end{remark}

\section{Orthogonal polynomials}

Let $I$ be an interval, $w$ a positive continuous function on its interior with $\int_I|t|^kw(t)\dd t<\infty$ for every
$k$, and $\langle f,g\rangle=\int_If(t)g(t)w(t)\dd t$ on real polynomials. This is an inner product: a non-zero
polynomial has finitely many zeros, so $\int f^2w>0$. Gram--Schmidt applied to $1,t,t^2,\dots$ produces, after rescaling
to make the leading coefficients $1$, the \emph{monic orthogonal polynomials} $p_0=1,p_1,p_2,\dots$: $p_n$ has degree $n$
and is orthogonal to every polynomial of degree less than $n$. They are unique, by Theorem~\ref{linear:thm:gramschmidt}.

\begin{example}[Legendre polynomials]\label{linear:exa:legendre}
For $I=[-1,1]$ and $w=1$, the polynomials
\[
P_n(t)=\frac1{2^nn!}\,\frac{\mathrm d^n}{\mathrm dt^n}(t^2-1)^n
\]
(Rodrigues, 1816) are orthogonal: for $k<n$, integrating by parts $k$ times,
\[
\int_{-1}^1t^k\frac{\mathrm d^n}{\mathrm dt^n}(t^2-1)^n\dd t=(-1)^kk!\int_{-1}^1\frac{\mathrm d^{n-k}}{\mathrm dt^{n-k}}(t^2-1)^n\dd t=0,
\]
because every derivative of $(t^2-1)^n=(t-1)^n(t+1)^n$ of order less than $n$ vanishes at $\pm1$, so all boundary terms
vanish, and the last integral is a difference of two such derivatives at $\pm1$. The Legendre polynomials are scaled so
that $P_n(1)=1$ (Exercise~\ref{linear:ex:legendrenorm}); the monic ones are $P_n$ divided by the leading coefficient
$\binom{2n}n/2^n$. Other weights give other famous families: $w=(1-t^2)^{-1/2}$ on $[-1,1]$ gives Chebyshev
(Exercise~\ref{linear:ex:chebyshev}), $e^{-t}$ on $[0,\infty)$ Laguerre, $e^{-t^2}$ on $\R$ Hermite.
\end{example}

\begin{figure}[ht]
\centering
\fig{linear-ch04-legendre}
\caption{The Legendre polynomials $P_0,\dots,P_5$. Each $P_n$ has $n$ simple zeros in $(-1,1)$, and the zeros of
consecutive polynomials interlace.}
\label{linear:fig:legendre}
\end{figure}

The picture suggests two theorems. Here is the first, which explains why only three polynomials ever talk to each other.

\begin{theorem}[three-term recurrence]\label{linear:thm:recurrence}
There are real numbers $\alpha_n$ and $\beta_n>0$ such that, with $p_{-1}=0$,
\[
p_{n+1}(t)=(t-\alpha_n)p_n(t)-\beta_np_{n-1}(t),\qquad \alpha_n=\frac{\langle tp_n,p_n\rangle}{\|p_n\|^2},\quad
\beta_n=\frac{\|p_n\|^2}{\|p_{n-1}\|^2}\ (n\ge1).
\]
\end{theorem}
\begin{proof}
The polynomial $tp_n-p_{n+1}$ has degree at most $n$ (both are monic of degree $n+1$), so it is a combination
$\sum_{k\le n}c_kp_k$ with $c_k=\langle tp_n,p_k\rangle/\|p_k\|^2$. The key point is that multiplication by $t$ is
symmetric: $\langle tp_n,p_k\rangle=\langle p_n,tp_k\rangle$, which vanishes when $k+1<n$ since $tp_k$ has degree $k+1$.
So only $c_n=\alpha_n$ and $c_{n-1}$ survive, and $\langle p_n,tp_{n-1}\rangle=\langle p_n,p_n\rangle$ because
$tp_{n-1}-p_n$ has degree less than $n$. Hence $c_{n-1}=\|p_n\|^2/\|p_{n-1}\|^2=\beta_n$.
\end{proof}

\begin{theorem}[zeros]\label{linear:thm:zeros}
The polynomial $p_n$ has $n$ distinct real zeros, all in the interior of $I$.
\end{theorem}
\begin{proof}
Let $t_1<\dots<t_k$ be the points of the interior of $I$ where $p_n$ changes sign (zeros of odd multiplicity), and
$q=\prod_{i=1}^k(t-t_i)$, with $q=1$ if $k=0$. Then $p_nq$ does not change sign on $I$ and is not zero, so
$\langle p_n,q\rangle=\int_Ip_nqw\ne0$. If $k<n$, $q$ has degree less than $n$ and $\langle p_n,q\rangle=0$. Hence $k\ge n$;
as $p_n$ has at most $n$ zeros, $k=n$, and these zeros are simple.
\end{proof}

There is a second proof, which explains the interlacing in Figure~\ref{linear:fig:legendre} and is also how the zeros are
computed in practice.

\begin{proposition}[the Jacobi matrix]\label{linear:prop:jacobi}
The zeros of $p_n$ are the eigenvalues of the real symmetric tridiagonal matrix
\[
J_n=\begin{pmatrix}\alpha_0&\sqrt{\beta_1}&&\\\sqrt{\beta_1}&\alpha_1&\ddots&\\&\ddots&\ddots&\sqrt{\beta_{n-1}}\\
&&\sqrt{\beta_{n-1}}&\alpha_{n-1}\end{pmatrix}.
\]
\end{proposition}
\begin{proof}
Let $T$ be the tridiagonal matrix with diagonal $\alpha_0,\dots,\alpha_{n-1}$, ones above the diagonal and
$\beta_1,\dots,\beta_{n-1}$ below it. The recurrence $tp_k=p_{k+1}+\alpha_kp_k+\beta_kp_{k-1}$ says that the vector
$v(t)=(p_0(t),\dots,p_{n-1}(t))^{\mathsf T}$ satisfies $Tv(t)=tv(t)-p_n(t)e_n$. At a zero $t$ of $p_n$, $v(t)$ is an
eigenvector of $T$ for the eigenvalue $t$ (it is non-zero since $p_0=1$). The $n$ distinct zeros (Theorem~\ref{linear:thm:zeros})
give $n$ distinct eigenvalues, so $\chi_T=p_n$. Finally $T=DJ_nD^{-1}$ with $D=\operatorname{diag}(d_0,\dots,d_{n-1})$,
$d_0=1$, $d_k=\sqrt{\beta_1\cdots\beta_k}$, as a direct check of the off-diagonal entries shows.
\end{proof}

Since $J_n$ is the top-left corner of $J_{n+1}$, the interlacing of the zeros of $p_n$ and $p_{n+1}$ is a special case of
Cauchy's interlacing theorem, proved in Appendix~E.

\section{Gauss quadrature}

To approximate $\int_If\,w$, a \emph{quadrature rule} chooses nodes $t_1,\dots,t_n$ and weights $\lambda_1,\dots,\lambda_n$
and uses $\sum_i\lambda_if(t_i)$. With $n$ nodes chosen in advance, interpolation makes the rule exact for polynomials
of degree less than $n$: take $\lambda_i=\int_I\ell_iw$, where $\ell_i$ is the Lagrange basis polynomial of
Example~\ref{linear:ex:lagrange}, $\ell_i(t_j)=\delta_{ij}$. Gauss's discovery (1814) is that the right choice of nodes doubles
the degree for free.

\begin{theorem}[Gauss quadrature]\label{linear:thm:gaussquad}
Let $t_1,\dots,t_n$ be the zeros of $p_n$ and $\lambda_i=\int_I\ell_iw$. Then
\[
\int_If(t)w(t)\dd t=\sum_{i=1}^n\lambda_if(t_i)\quad\text{for every polynomial $f$ of degree at most $2n-1$.}
\]
The weights are positive, $\lambda_i=\int_I\ell_i^2w$, and no rule with $n$ nodes is exact for all polynomials of degree
$2n$.
\end{theorem}
\begin{proof}
Divide: $f=qp_n+r$ with $\deg q\le n-1$ and $\deg r\le n-1$. Then $\int fw=\langle q,p_n\rangle+\int rw=\int rw$, by
orthogonality. Since $f(t_i)=r(t_i)$ and the rule is exact for $r$ (interpolation), $\sum\lambda_if(t_i)=\sum\lambda_ir(t_i)
=\int rw=\int fw$. Applying this to $\ell_i^2$, of degree $2n-2$, gives $\int\ell_i^2w=\sum_j\lambda_j\ell_i(t_j)^2=\lambda_i$,
which is positive. Finally, for any nodes $s_1,\dots,s_n$, the polynomial $\prod_i(t-s_i)^2$ of degree $2n$ has a positive
integral and vanishes at every node.
\end{proof}

With $n=2$ and Legendre's weight, $p_2=t^2-\frac13$, the nodes are $\pm\frac1{\sqrt3}$ and the weights are $1$: the rule
$f(-\frac1{\sqrt3})+f(\frac1{\sqrt3})$ integrates every cubic over $[-1,1]$ exactly, with two evaluations. Positivity of the
weights matters in practice: it makes the rule stable, and it gives convergence for every continuous $f$ (Exercise~\ref{linear:ex:gaussconv}).

\section{Adjoints and isometries}\label{linear:sec:adjoint}

From now on $V$ is a finite-dimensional inner product space.

\begin{theorem}[Riesz representation]\label{linear:thm:riesz}
For every linear form $\varphi$ on $V$ there is a unique $y\in V$ with $\varphi(x)=\langle x,y\rangle$ for all $x$.
\end{theorem}
\begin{proof}
With an orthonormal basis $(e_i)$, set $y=\sum_i\overline{\varphi(e_i)}e_i$. For $x=\sum_ix_ie_i$,
$\langle x,y\rangle=\sum_ix_i\varphi(e_i)=\varphi(x)$. If $y'$ also works, $\langle x,y-y'\rangle=0$ for all $x$, including
$x=y-y'$, so $y=y'$.
\end{proof}

So the inner product identifies $V$ with its dual $V^*$ (conjugate-linearly when $\mathbb K=\C$), with no choice of basis.
For an operator $A$ on $V$, the form $x\mapsto\langle Ax,y\rangle$ is represented by a unique vector, called $A^*y$.

\begin{proposition}[the adjoint]\label{linear:prop:adjoint}
There is a unique linear map $A^*$ with $\langle Ax,y\rangle=\langle x,A^*y\rangle$ for all $x,y$. In an orthonormal
basis its matrix is the conjugate transpose of that of $A$. Moreover $(AB)^*=B^*A^*$, $A^{**}=A$,
\[
\ker A^*=(\operatorname{Im}A)^\perp\quad\text{and}\quad\operatorname{Im}A^*=(\ker A)^\perp.
\]
\end{proposition}
\begin{proof}
Existence and uniqueness come from Theorem~\ref{linear:thm:riesz}, and linearity of $A^*$ from uniqueness. In an orthonormal
basis the $(i,j)$ entry of $A^*$ is $\langle A^*e_j,e_i\rangle=\overline{\langle e_i,A^*e_j\rangle}=\overline{\langle
Ae_i,e_j\rangle}=\overline{a_{ji}}$. The two product rules follow from the defining identity. Finally, $y\in\ker A^*$ if
and only if $\langle x,A^*y\rangle=\langle Ax,y\rangle=0$ for all $x$, that is $y\perp\operatorname{Im}A$. Applying this to
$A^*$ and taking orthogonals gives the last identity, as $W^{\perp\perp}=W$ in finite dimension.
\end{proof}

This is Theorem~\ref{linear:thm:transpose} of Chapter~\ref{linear:ch:duality} once the dual is identified with $V$; the Fredholm
alternative becomes ``$Ax=b$ is solvable if and only if $b$ is orthogonal to every solution of $A^*y=0$''.

An operator $U$ is \emph{unitary} (\emph{orthogonal} when $\mathbb K=\R$) if it preserves inner products:
$\langle Ux,Uy\rangle=\langle x,y\rangle$.

\begin{theorem}[unitary operators]\label{linear:thm:unitary}
For an operator $U$ on $V$, the following are equivalent: (i) $U$ preserves inner products; (ii) $U$ preserves norms;
(iii) $U$ maps some (equivalently every) orthonormal basis to an orthonormal basis; (iv) $U^*U=I$.
\end{theorem}
\begin{proof}
(i)$\Rightarrow$(ii) is the case $x=y$. (ii)$\Rightarrow$(i) by the \emph{polarisation identities}, which recover the inner
product from the norm:
\[
\operatorname{Re}\langle x,y\rangle=\tfrac14\bigl(\|x+y\|^2-\|x-y\|^2\bigr),\qquad \operatorname{Im}\langle x,y\rangle=
\operatorname{Re}\langle x,iy\rangle.
\]
(i)$\Rightarrow$(iii) is clear. (iii)$\Rightarrow$(i): if $(Ue_i)$ is orthonormal, then for $x=\sum x_ie_i$ and $y=\sum
y_ie_i$, both $\langle Ux,Uy\rangle$ and $\langle x,y\rangle$ equal $\sum_ix_i\bar y_i$. (i)$\Leftrightarrow$(iv):
$\langle Ux,Uy\rangle=\langle U^*Ux,y\rangle$, and this equals $\langle x,y\rangle$ for all $x,y$ if and only if $U^*U=I$.
\end{proof}

Unitary matrices are those with orthonormal columns; they form a group, and every change of orthonormal basis is one.
Requiring linearity turns out to be unnecessary in the real case.

\begin{theorem}\label{linear:thm:isometry}
Let $f:\R^n\to\R^n$ be any map with $f(0)=0$ and $\|f(x)-f(y)\|=\|x-y\|$ for all $x,y$. Then $f$ is linear, hence
orthogonal.
\end{theorem}
\begin{proof}
Taking $y=0$, $f$ preserves norms. Since $\langle x,y\rangle=\frac12\bigl(\|x\|^2+\|y\|^2-\|x-y\|^2\bigr)$ in a real space,
$f$ preserves inner products. Now expand $\|f(x+y)-f(x)-f(y)\|^2$ into nine inner products of the vectors
$f(x+y),f(x),f(y)$; each equals the corresponding inner product of $x+y,x,y$, so the total equals
$\|(x+y)-x-y\|^2=0$. Hence $f(x+y)=f(x)+f(y)$, and in the same way $\|f(\lambda x)-\lambda f(x)\|^2=\|\lambda x-\lambda
x\|^2=0$.
\end{proof}

Thus a rigid motion of space is a translation followed by an orthogonal map: the starting point of the geometry of
Tome~8. To close the section, here is what Gram--Schmidt says about inner products themselves. On $\mathbb K^n$, every
inner product has the form $\langle x,y\rangle_G=y^*Gx$, where $G=(\langle e_j,e_i\rangle_G)$ is the Gram matrix of the
standard basis: a Hermitian matrix ($G^*=G$) with $x^*Gx>0$ for $x\ne0$, that is, \emph{positive definite}. Conversely
every positive definite $G$ defines an inner product this way.

\begin{corollary}[Cholesky factorisation]\label{linear:cor:cholesky}
Every positive definite matrix $G$ is $G=R^*R$ for a unique upper triangular $R$ with positive diagonal.
\end{corollary}
\begin{proof}
Apply Gram--Schmidt to the standard basis $e_1,\dots,e_n$ for the inner product $\langle\cdot,\cdot\rangle_G$, obtaining
an orthonormal basis $f_1,\dots,f_n$ with $e_j=\sum_{i\le j}r_{ij}f_i$, $r_{jj}>0$. With $F=(f_1|\cdots|f_n)$, this says
$I=FR$, so $F=R^{-1}$; and orthonormality says $F^*GF=I$. Hence $G=(F^*)^{-1}F^{-1}=R^*R$. Uniqueness: if $G=R^*R$, then
the columns of $R^{-1}$ are $G$-orthonormal and obtained from the standard basis by an upper triangular change of
basis with positive diagonal, so they are the Gram--Schmidt basis by the uniqueness in Theorem~\ref{linear:thm:gramschmidt}.
\end{proof}

Appendix~K returns to positive matrices, with Hadamard's and Minkowski's inequalities, and uses Cholesky constantly.

\begin{beauty}
Gram--Schmidt is a dull algorithm, and yet: run it on the monomials $1,t,t^2,\dots$ and out come the Legendre
polynomials, whose zeros are the eigenvalues of a tridiagonal matrix, interlace like the eigenvalues of nested
matrices, and are exactly the points where Gauss's rule integrates polynomials of twice the expected degree. The volumes
of the parallelepipeds built on those same monomials are the Cauchy determinants of the Hilbert matrix, rational numbers
whose inverses are integers. One right angle, dropped over and over, organises all of it.
\end{beauty}

\section{Exercises}

\begin{exercise}[$\star$]\label{linear:ex:polar}
(a) Prove the parallelogram law $\|x+y\|^2+\|x-y\|^2=2\|x\|^2+2\|y\|^2$ in any inner product space, and interpret it for
a parallelogram in the plane. (b) In a complex space, show that $\langle x,y\rangle=\frac14\sum_{k=0}^3i^k\|x+i^ky\|^2$.
\end{exercise}

\begin{exercise}[$\star\star\star$]\label{linear:ex:jvn}
(Jordan--von Neumann.) Let $\|\cdot\|$ be a norm on a real vector space satisfying the parallelogram law. Show that
$\langle x,y\rangle=\frac14(\|x+y\|^2-\|x-y\|^2)$ is an inner product whose norm is $\|\cdot\|$.
\end{exercise}

\begin{exercise}[$\star$]\label{linear:ex:qr}
Compute the $QR$ factorisation of
$A=\begin{pmatrix}1&1&0\\1&0&1\\0&1&1\end{pmatrix}$, and check that $|\det A|=r_{11}r_{22}r_{33}$.
\end{exercise}

\begin{exercise}[$\star$]\label{linear:ex:projmatrix}
Let $A\in M_{m,n}(\R)$ have independent columns. Show that $P=A(A^{\mathsf T}A)^{-1}A^{\mathsf T}$ is the matrix of the
orthogonal projection onto $\operatorname{Im}A$. Conversely, show that a matrix with $P^2=P=P^{\mathsf T}$ is the
orthogonal projection onto its image, and that $\operatorname{tr}P=\operatorname{rank}P$.
\end{exercise}

\begin{exercise}[$\star\star$]\label{linear:ex:parabola}
Fit $y=a+bx+cx^2$ by least squares to the five points $(-2,3)$, $(-1,1)$, $(0,0)$, $(1,1)$, $(2,4)$. Do it a second time
using the basis $1$, $x$, $x^2-2$, and explain why the normal equations become diagonal.
\end{exercise}

\begin{exercise}[$\star$]\label{linear:ex:distpoly}
Find the polynomial $a+bt$ minimising $\int_0^1(t^2-a-bt)^2\dd t$, and check the minimum against
Corollary~\ref{linear:cor:hilbert}.
\end{exercise}

\begin{exercise}[$\star$]\label{linear:ex:hilbertint}
Show that $\det H_{n+1}/\det H_n=\dfrac1{(2n+1)\binom{2n}n^2}$ and deduce that $1/\det H_n$ is an integer. Compute it for
$n\le4$.
\end{exercise}

\begin{exercise}[$\star\star$]\label{linear:ex:obtuse}
Show that $\R^n$ contains $n+1$ vectors with pairwise negative inner products, but not $n+2$. (Project onto the
orthogonal of one of them.)
\end{exercise}

\begin{exercise}[$\star\star$]\label{linear:ex:basel}
In the space of continuous real functions on $[-\pi,\pi]$ with $\langle f,g\rangle=\int_{-\pi}^\pi fg$, show that the
functions $\frac{\sin kt}{\sqrt\pi}$, $k\ge1$, are orthonormal. Apply Bessel's inequality~\eqref{linear:eq:bessel} to $f(t)=t$ to
prove $\sum_{k\ge1}\frac1{k^2}\le\frac{\pi^2}6$. (Equality is Parseval's identity, Tome~7.)
\end{exercise}

\begin{exercise}[$\star\star$]\label{linear:ex:legendrenorm}
Show that $P_n(1)=1$, that the leading coefficient of $P_n$ is $\binom{2n}n/2^n$, and that $\int_{-1}^1P_n^2=\frac2{2n+1}$.
(Use $\int_{-1}^1(1-t^2)^n\dd t=\frac{2^{2n+1}(n!)^2}{(2n+1)!}$, which you may prove by induction.)
\end{exercise}

\begin{exercise}[$\star\star$]\label{linear:ex:chebyshev}
Let $T_n(\cos\theta)=\cos n\theta$. Show that $T_n$ is a polynomial of degree $n$, that the $T_n$ are orthogonal for the
weight $(1-t^2)^{-1/2}$ on $(-1,1)$, and that the Gauss rule for this weight has nodes $\cos\frac{(2k-1)\pi}{2n}$ and all
weights equal to $\frac\pi n$.
\end{exercise}

\begin{exercise}[$\star\star$]\label{linear:ex:gaussconv}
Let $G_n(f)=\sum_i\lambda_i^{(n)}f(t_i^{(n)})$ be the $n$-point Gauss rule on a bounded interval $I$. Using positivity of the
weights and the Weierstrass approximation theorem, show that $G_n(f)\to\int_Ifw$ for every continuous $f$ on $I$.
\end{exercise}

\begin{exercise}[$\star\star$]\label{linear:ex:householder}
For $u\ne0$ in $\R^n$, let $H_u=I-2\frac{uu^{\mathsf T}}{u^{\mathsf T}u}$. Show that $H_u$ is orthogonal and symmetric (a
reflection), that $H_{x-y}x=y$ when $\|x\|=\|y\|$ and $x\ne y$, and that every orthogonal $n\times n$ matrix is a product of
at most $n$ reflections.
\end{exercise}

\begin{exercise}[$\star\star\star$]\label{linear:ex:muntz}
Let $0\le a_1<a_2<\dots$ be real and $m\ge0$ an integer distinct from all $a_k$. In $L^2[0,1]$, show that the distance
$d_N$ from $t^m$ to $\operatorname{span}(t^{a_1},\dots,t^{a_N})$ satisfies
\[
d_N^2=\frac1{2m+1}\prod_{k=1}^N\Bigl(\frac{a_k-m}{a_k+m+1}\Bigr)^2.
\]
Deduce that if $a_k\to\infty$ and $\sum_k\frac1{a_k}=\infty$, then $d_N\to0$ (one half of M\"untz's theorem, 1914).
\end{exercise}

\section{Hints}

\begin{hint}{linear:ex:polar}
Expand both squares and add. For (b), expand each $\|x+i^ky\|^2$ and use $\sum_ki^k=0$.
\end{hint}
\begin{hint}{linear:ex:jvn}
Show $\langle x+z,y\rangle+\langle x-z,y\rangle=2\langle x,y\rangle$, then $\langle2x,y\rangle=2\langle x,y\rangle$, then
additivity, then homogeneity for rational and finally real scalars by continuity.
\end{hint}
\begin{hint}{linear:ex:qr}
Gram--Schmidt on the columns; $\det Q=\pm1$.
\end{hint}
\begin{hint}{linear:ex:projmatrix}
$Pb=Ax$ for the least squares solution $x$. For the converse, check $b-Pb\perp\operatorname{Im}P$; for the trace, use a basis
adapted to $\operatorname{Im}P\oplus\ker P$.
\end{hint}
\begin{hint}{linear:ex:parabola}
On these five points, $\sum x=\sum x^3=0$ and $\sum(x^2-2)=0$.
\end{hint}
\begin{hint}{linear:ex:distpoly}
The error $t^2-a-bt$ must be orthogonal to $1$ and to $t$.
\end{hint}
\begin{hint}{linear:ex:hilbertint}
The ratio was computed in the proof of Corollary~\ref{linear:cor:hilbert}.
\end{hint}
\begin{hint}{linear:ex:obtuse}
If $v_1,\dots,v_m$ are pairwise obtuse, the projections of $v_1,\dots,v_{m-1}$ onto $v_m^\perp$ are still pairwise obtuse.
\end{hint}
\begin{hint}{linear:ex:basel}
$\int_{-\pi}^\pi t\sin kt\dd t=\frac{2\pi(-1)^{k+1}}k$, by parts.
\end{hint}
\begin{hint}{linear:ex:legendrenorm}
Leibniz's rule on $(t-1)^n(t+1)^n$ for $P_n(1)$. For the norm, $\int P_n^2$ equals the leading coefficient times
$\int t^nP_n$.
\end{hint}
\begin{hint}{linear:ex:chebyshev}
$\cos(n+1)\theta+\cos(n-1)\theta=2\cos\theta\cos n\theta$; substitute $t=\cos\theta$ in the integrals.
\end{hint}
\begin{hint}{linear:ex:gaussconv}
$\sum_i\lambda_i^{(n)}=\int_Iw$ for all $n$. Approximate $f$ uniformly within $\varepsilon$ by a polynomial $q$, and compare
$G_n(f)$ with $G_n(q)$ once $2n-1\ge\deg q$.
\end{hint}
\begin{hint}{linear:ex:householder}
For the last part, induct on $n$: send $Qe_1$ back to $e_1$ by one reflection.
\end{hint}
\begin{hint}{linear:ex:muntz}
Theorem~\ref{linear:thm:gram}(b) with Cauchy's determinant: $\langle t^a,t^b\rangle=\frac1{a+b+1}$. For the limit, write each factor
as $1-\frac{2m+1}{a_k+m+1}$ and use $\log(1-u)\le-u$.
\end{hint}

\section{Solutions}

\begin{solution}{linear:ex:polar}
\idea Every identity about norms in an inner product space is proved by expanding $\|u\|^2=\langle u,u\rangle$.
\writeup (a) $\|x\pm y\|^2=\|x\|^2\pm2\operatorname{Re}\langle x,y\rangle+\|y\|^2$; adding cancels the middle terms. For a
parallelogram with sides $x,y$, the diagonals are $x+y$ and $x-y$: the sum of the squares of the diagonals equals the sum
of the squares of the four sides.

(b) $\|x+i^ky\|^2=\|x\|^2+\|y\|^2+\overline{i^k}\langle x,y\rangle+i^k\langle y,x\rangle$. Multiply by $i^k$ and sum over
$k=0,1,2,3$: the terms $i^k(\|x\|^2+\|y\|^2)$ and $i^{2k}\langle y,x\rangle$ sum to zero (as $\sum i^k=\sum(-1)^k=0$),
while $i^k\overline{i^k}=1$ leaves $4\langle x,y\rangle$.
\end{solution}

\begin{solution}{linear:ex:jvn}
\idea The formula is forced by polarisation; symmetry and $\langle x,x\rangle=\|x\|^2$ are immediate, so all the work is
linearity in $x$. The parallelogram law is the only tool, so apply it to the pairs that occur in
$\|x\pm z+y\|$ and $\|x\pm z-y\|$.
\writeup Symmetry is clear, $\langle x,x\rangle=\frac14\|2x\|^2=\|x\|^2$, and $\langle0,y\rangle=0$. The parallelogram
law for the vectors $x+y$ and $z$ gives $\|x+z+y\|^2+\|x-z+y\|^2=2\|x+y\|^2+2\|z\|^2$; for $x-y$ and $z$ it gives
$\|x+z-y\|^2+\|x-z-y\|^2=2\|x-y\|^2+2\|z\|^2$. Subtracting and dividing by $4$,
\[
\langle x+z,y\rangle+\langle x-z,y\rangle=2\langle x,y\rangle.
\]
With $z=x$ this gives $\langle2x,y\rangle=2\langle x,y\rangle$. For arbitrary $u,v$, put $x=\frac{u+v}2$, $z=\frac{u-v}2$:
then $\langle u,y\rangle+\langle v,y\rangle=2\langle\frac{u+v}2,y\rangle=\langle u+v,y\rangle$, so $\langle\cdot,y\rangle$
is additive. Additivity gives $\langle kx,y\rangle=k\langle x,y\rangle$ for integers $k$ (using $\langle-x,y\rangle=-\langle
x,y\rangle$, clear from the formula), then for rationals ($q\langle\frac pqx,y\rangle=\langle px,y\rangle$). Finally
$\lambda\mapsto\langle\lambda x,y\rangle$ is continuous, since the norm is continuous ($|\|u\|-\|v\||\le\|u-v\|$), and agrees
with $\lambda\langle x,y\rangle$ on $\Q$, hence on $\R$. Positivity is $\langle x,x\rangle=\|x\|^2>0$ for $x\ne0$.
\end{solution}

\begin{solution}{linear:ex:qr}
\idea Run Gram--Schmidt and record the coefficients: $r_{ij}=\langle a_j,e_i\rangle$.
\writeup $e_1=\frac1{\sqrt2}(1,1,0)$. Then $a_2-\frac12(1,1,0)=(\frac12,-\frac12,1)$, of norm $\sqrt{3/2}$, so
$e_2=\frac1{\sqrt6}(1,-1,2)$. Then $\langle a_3,e_1\rangle=\frac1{\sqrt2}$, $\langle a_3,e_2\rangle=\frac1{\sqrt6}$, and
$a_3-\frac12(1,1,0)-\frac16(1,-1,2)=\frac23(-1,1,1)$, of norm $\frac2{\sqrt3}$, so $e_3=\frac1{\sqrt3}(-1,1,1)$. Hence
\[
Q=\begin{pmatrix}\frac1{\sqrt2}&\frac1{\sqrt6}&-\frac1{\sqrt3}\\[2pt]\frac1{\sqrt2}&-\frac1{\sqrt6}&\frac1{\sqrt3}\\[2pt]0&\frac2{\sqrt6}&\frac1{\sqrt3}\end{pmatrix},\qquad
R=\begin{pmatrix}\sqrt2&\frac1{\sqrt2}&\frac1{\sqrt2}\\[2pt]0&\sqrt{\frac32}&\frac1{\sqrt6}\\[2pt]0&0&\frac2{\sqrt3}\end{pmatrix}.
\]
Then $r_{11}r_{22}r_{33}=\sqrt2\cdot\sqrt{3/2}\cdot\frac2{\sqrt3}=2$, and $\det A=1\cdot(0-1)-1\cdot(1-0)=-2$. In general
$|\det A|=|\det Q|\det R=\prod r_{jj}$, the volume as base times heights; here $\det Q=-1$.
\end{solution}

\begin{solution}{linear:ex:projmatrix}
\idea A projection is determined by what it does: fix the subspace, kill its orthogonal.
\writeup By Theorem~\ref{linear:thm:lsq}, the closest point of $\operatorname{Im}A$ to $b$ is $Ax$ with $x=(A^{\mathsf T}A)^{-1}
A^{\mathsf T}b$, which is $Pb$. Conversely, let $P^2=P=P^{\mathsf T}$ and $W=\operatorname{Im}P$. For $w=Pv\in W$,
$Pw=P^2v=w$. For any $b$ and $w\in W$, $\langle b-Pb,w\rangle=\langle b-Pb,Pw\rangle=\langle P(b-Pb),w\rangle=\langle
Pb-P^2b,w\rangle=0$. So $Pb\in W$ and $b-Pb\perp W$, which characterises $P_Wb$. In a basis made of a basis of $W$ and a
basis of $W^\perp$, $P=\operatorname{diag}(1,\dots,1,0,\dots,0)$ with $\dim W$ ones, so $\operatorname{tr}P=\dim W=
\operatorname{rank}P$.
\end{solution}

\begin{solution}{linear:ex:parabola}
\idea Normal equations in the monomial basis first; then notice that in an orthogonal basis the coefficients are just
projections, one inner product each.
\writeup With the sums $\sum1=5$, $\sum x=\sum x^3=0$, $\sum x^2=10$, $\sum x^4=34$, $\sum y=9$, $\sum xy=2$, $\sum
x^2y=30$, the normal equations are $5a+10c=9$, $10b=2$, $10a+34c=30$, so $b=\frac15$, $c=\frac67$, $a=\frac3{35}$.

The vectors of values of $1$, $x$, $x^2-2$ at the five points are pairwise orthogonal ($\sum x=0$, $\sum(x^2-2)=0$,
$\sum x(x^2-2)=0$), so the projection of $y$ onto their span is
\[
\frac{\langle y,1\rangle}{5}+\frac{\langle y,x\rangle}{10}x+\frac{\langle y,x^2-2\rangle}{14}(x^2-2)=\frac95+\frac x5+
\frac67(x^2-2),
\]
which is $\frac3{35}+\frac x5+\frac67x^2$ again. The Gram matrix of an orthogonal basis is diagonal, which is why discrete
orthogonal polynomials were how least squares fits were computed by hand for a century.
\end{solution}

\begin{solution}{linear:ex:distpoly}
\idea Orthogonality of the error to $1$ and $t$ gives two linear equations; Corollary~\ref{linear:cor:hilbert} predicts the
answer $\frac1{6\sqrt5}$ for the distance.
\writeup $\int_0^1(t^2-a-bt)\dd t=\frac13-a-\frac b2=0$ and $\int_0^1(t^2-a-bt)t\dd t=\frac14-\frac a2-\frac b3=0$ give
$b=1$, $a=-\frac16$. The error is $t^2-t+\frac16$, which with $s=t-\frac12$ is $s^2-\frac1{12}$, and
\[
\int_{-1/2}^{1/2}\Bigl(s^2-\frac1{12}\Bigr)^2\dd s=\frac1{80}-\frac2{12}\cdot\frac1{12}+\frac1{144}=\frac1{80}-\frac1{144}=\frac1{180}.
\]
The corollary with $n=2$ gives $\bigl(\binom42\sqrt5\bigr)^{-2}=\frac1{180}$.
\end{solution}

\begin{solution}{linear:ex:hilbertint}
\idea The ratio is the squared distance of Corollary~\ref{linear:cor:hilbert}; rewrite $\frac{(n!)^4}{(2n)!(2n+1)!}$ with a
binomial coefficient.
\writeup $\frac{(n!)^4}{(2n)!(2n+1)!}=\frac{(n!)^4}{(2n)!^2(2n+1)}=\frac1{(2n+1)\binom{2n}n^2}$. Since $\det H_1=1$,
\[
\frac1{\det H_n}=\prod_{k=1}^{n-1}(2k+1)\binom{2k}k^2,
\]
a product of integers. The values for $n=1,2,3,4$ are $1$, $3\cdot4=12$, $12\cdot5\cdot36=2160$ and $2160\cdot7\cdot400=
6048000$.
\end{solution}

\begin{solution}{linear:ex:obtuse}
\idea For the construction, think of the vertices of a regular simplex seen from its centre. For the bound, removing
one vector by projection should keep the others obtuse, which sets up an induction on the dimension.
\writeup \emph{Construction.} Let $f_1,\dots,f_{n+1}$ be the standard basis of $\R^{n+1}$ and $v_i=f_i-\frac1{n+1}
\sum_jf_j$. These lie in the hyperplane $\sum x_j=0$, a copy of $\R^n$, and $\langle v_i,v_j\rangle=-\frac1{n+1}<0$ for
$i\ne j$.

\emph{Bound}, by induction on $n$. In $\R^1$, three non-zero numbers cannot have pairwise negative products (two of them
have the same sign). Let $v_1,\dots,v_m\in\R^n$ be pairwise obtuse, and let $w_i=v_i-\frac{\langle v_i,v_m\rangle}
{\|v_m\|^2}v_m$ for $i<m$, the projections onto $v_m^\perp\cong\R^{n-1}$. Then
\[
\langle w_i,w_j\rangle=\langle v_i,v_j\rangle-\frac{\langle v_i,v_m\rangle\langle v_j,v_m\rangle}{\|v_m\|^2}<0,
\]
as the first term is negative and the subtracted product of two negative numbers is positive. In particular the $w_i$
are non-zero, and by induction $m-1\le(n-1)+1$, so $m\le n+1$.
\end{solution}

\begin{solution}{linear:ex:basel}
\idea Bessel's inequality turns any orthonormal family into a lower bound for $\int f^2$; with $f(t)=t$ and the sines,
the coefficients are $\frac1k$ up to constants, so out comes $\sum\frac1{k^2}$.
\writeup $\int_{-\pi}^\pi\sin kt\sin lt\dd t=\frac12\int_{-\pi}^\pi(\cos(k-l)t-\cos(k+l)t)\dd t=\pi\delta_{kl}$ for $k,l\ge1$,
so the $\frac{\sin kt}{\sqrt\pi}$ are orthonormal. By parts, $\int_{-\pi}^\pi t\sin kt\dd t=\bigl[-\frac{t\cos kt}k\bigr]
_{-\pi}^\pi+\int\frac{\cos kt}k=\frac{2\pi(-1)^{k+1}}k$. Bessel's inequality gives
\[
\sum_{k\ge1}\frac1\pi\cdot\frac{4\pi^2}{k^2}\le\int_{-\pi}^\pi t^2\dd t=\frac{2\pi^3}3,\qquad\text{that is}\qquad
\sum_{k\ge1}\frac1{k^2}\le\frac{\pi^2}6.
\]
\end{solution}

\begin{solution}{linear:ex:legendrenorm}
\idea Evaluating an $n$-th derivative of a product at a root of one factor kills almost every Leibniz term. For the
norm, orthogonality lets us replace one copy of $P_n$ by its leading term.
\writeup By Leibniz, $\frac{\mathrm d^n}{\mathrm dt^n}\bigl((t-1)^n(t+1)^n\bigr)=\sum_k\binom nk\frac{\mathrm d^k}{\mathrm
dt^k}(t-1)^n\cdot\frac{\mathrm d^{n-k}}{\mathrm dt^{n-k}}(t+1)^n$; at $t=1$ only $k=n$ survives, giving $n!\,2^n$, so
$P_n(1)=1$. The leading term of $(t^2-1)^n$ is $t^{2n}$, whose $n$-th derivative is $\frac{(2n)!}{n!}t^n$; dividing by
$2^nn!$ gives the leading coefficient $\binom{2n}n/2^n$. Since $P_n\perp$ lower degrees, $\int P_n^2=\frac{\binom{2n}n}{2^n}
\int t^nP_n$, and by the computation of Example~\ref{linear:exa:legendre} with $k=n$,
\[
\int_{-1}^1t^nP_n=\frac{(-1)^nn!}{2^nn!}\int_{-1}^1(t^2-1)^n=\frac1{2^n}\int_{-1}^1(1-t^2)^n=\frac{2^{n+1}(n!)^2}{(2n+1)!}.
\]
Multiplying, $\int P_n^2=\frac{(2n)!}{2^n(n!)^2}\cdot\frac{2^{n+1}(n!)^2}{(2n+1)!}=\frac2{2n+1}$. The integral $\int_{-1}^1(1-t^2)^n$
follows by induction from the integration by parts $I_n=\frac{2n}{2n+1}I_{n-1}$.
\end{solution}

\begin{solution}{linear:ex:chebyshev}
\idea The substitution $t=\cos\theta$ turns the weight $(1-t^2)^{-1/2}\dd t$ into $\dd\theta$, and everything becomes
the orthogonality of cosines.
\writeup From $T_{n+1}+T_{n-1}=2tT_n$, $T_0=1$ and $T_1=t$, induction shows $T_n$ is a polynomial of degree $n$
(leading coefficient $2^{n-1}$ for $n\ge1$). With $t=\cos\theta$,
\[
\int_{-1}^1T_mT_n\frac{\dd t}{\sqrt{1-t^2}}=\int_0^\pi\cos m\theta\cos n\theta\dd\theta=0\quad(m\ne n).
\]
So $p_n=2^{1-n}T_n$, whose zeros are the $t$ with $\cos n\theta=0$, that is $\theta=\frac{(2k-1)\pi}{2n}$, $k=1,\dots,n$.
For the weights, the rule must integrate $T_j$ exactly for $j\le2n-1$. With $\theta_k=\frac{(2k-1)\pi}{2n}$ and $0<j<2n$,
$\sum_k\cos j\theta_k=\operatorname{Re}\bigl(e^{ij\pi/2n}\sum_{k=0}^{n-1}e^{ij\pi k/n}\bigr)=\operatorname{Re}\frac{e^{ij\pi/2n}
(1-e^{ij\pi})}{1-e^{ij\pi/n}}=\operatorname{Re}\frac{1-(-1)^j}{e^{-ij\pi/2n}-e^{ij\pi/2n}}$, and the last fraction is purely
imaginary, so the sum is $0$. Thus equal weights $\frac\pi n$ integrate $T_0$ correctly ($\int\frac{\dd t}{\sqrt{1-t^2}}=\pi$)
and every $T_j$, $0<j<2n$, correctly (both sides $0$); by uniqueness of the weights of an interpolatory rule, these are the
Gauss weights.
\end{solution}

\begin{solution}{linear:ex:gaussconv}
\idea Exactness handles polynomials; positivity of weights says the rule cannot amplify errors, since it is an average.
\writeup Applying the rule to $f=1$ gives $\sum_i\lambda_i^{(n)}=\int_Iw=:M$. Given $\varepsilon>0$, Weierstrass gives a
polynomial $q$ with $|f-q|\le\varepsilon$ on $I$. For $2n-1\ge\deg q$, $G_n(q)=\int qw$, and since the weights are
positive,
\[
\Bigl|G_n(f)-\int_Ifw\Bigr|\le\sum_i\lambda_i^{(n)}|f-q|(t_i^{(n)})+\int_I|q-f|w\le2M\varepsilon.
\]
\end{solution}

\begin{solution}{linear:ex:householder}
\idea A reflection fixes a hyperplane and reverses its normal. To send $x$ to $y$, reflect in the perpendicular bisector,
whose normal is $x-y$.
\writeup $H_u$ is symmetric, and $H_u^2=I-4\frac{uu^{\mathsf T}}{u^{\mathsf T}u}+4\frac{u(u^{\mathsf T}u)u^{\mathsf T}}{(u^{\mathsf T}u)^2}=I$,
so $H_u^{\mathsf T}H_u=I$. With $u=x-y$ and $\|x\|=\|y\|$, $u^{\mathsf T}u=2\|x\|^2-2\langle x,y\rangle=2\langle x,u\rangle$, so
$H_ux=x-2\frac{\langle x,u\rangle}{u^{\mathsf T}u}u=x-u=y$. For the factorisation, induct on $n$ (the case $n=1$: $Q=\pm1$,
a product of at most one reflection). Given $Q$ orthogonal, if $Qe_1\ne e_1$ let $H=H_{Qe_1-e_1}$, so $HQe_1=e_1$;
otherwise $H=I$. Then $HQ$ is orthogonal and fixes $e_1$, hence preserves $e_1^\perp$, and acts there as an orthogonal
map of $\R^{n-1}$, a product of at most $n-1$ reflections of $e_1^\perp$ by induction; each extends to a reflection of
$\R^n$ fixing $e_1$. So $Q=H^{-1}(HQ)=H(HQ)$ is a product of at most $n$ reflections. This is the Cartan--Dieudonn\'e
theorem, and it is how Householder computes $QR$ factorisations stably.
\end{solution}

\begin{solution}{linear:ex:muntz}
\idea The distance is a ratio of Gram determinants, and here every Gram matrix is a Cauchy matrix. Adding one exponent
multiplies the Cauchy determinant by a simple factor, exactly as in the proof of Corollary~\ref{linear:cor:hilbert}.
\writeup $\langle t^a,t^b\rangle=\frac1{a+b+1}$, so the Gram matrix of $t^{a_1},\dots,t^{a_N}$ is Cauchy with $x_i=a_i$,
$y_j=a_j+1$, and adding the exponent $m$ as $(N+1)$-st multiplies its determinant by
\[
\frac{\prod_k(m-a_k)^2}{(2m+1)\prod_k(a_k+m+1)^2}
\]
by Theorem~\ref{linear:thm:cauchydet} (the factors involving the new index; order does not matter as everything is squared).
Theorem~\ref{linear:thm:gram}(b) gives the formula for $d_N^2$. Each factor is $\bigl(1-\frac{2m+1}{a_k+m+1}\bigr)^2$; for $k$
large it lies in $(0,1)$, and $\log(1-u)\le-u$ gives $\log d_N^2\le C-2\sum_{k\le N}\frac{2m+1}{a_k+m+1}\to-\infty$ when
$\sum\frac1{a_k}=\infty$ (as $\frac1{a_k+m+1}\ge\frac1{2a_k}$ once $a_k\ge m+1$). So $d_N\to0$. Combined with the density
of polynomials in $L^2[0,1]$, this shows the span of the $t^{a_k}$ is dense; M\"untz and Sz\'asz proved the converse.
\end{solution}

\section*{Sources and further reading}

\begin{itemize}
\item A.-M.~Legendre, \emph{Nouvelles m\'ethodes pour la d\'etermination des orbites des com\`etes}, Paris, 1805, appendix
``Sur la m\'ethode des moindres quarr\'es''.
\item C.~F.~Gauss, \emph{Theoria motus corporum coelestium}, 1809, and \emph{Theoria combinationis observationum erroribus
minimis obnoxiae}, 1821--1823 (the Gauss--Markov theorem); \emph{Methodus nova integralium valores per approximationem
inveniendi}, 1814 (Gauss quadrature).
\item S.~M.~Stigler, \emph{Gauss and the invention of least squares}, Annals of Statistics 9 (1981), 465--474; D.~Teets and
K.~Whitehead, \emph{The discovery of Ceres: how Gauss became famous}, Mathematics Magazine 72 (1999), 83--93.
\item A.-L.~Cauchy, \emph{M\'emoire sur les fonctions altern\'ees et sur les sommes altern\'ees}, Exercices d'analyse et de
physique math\'ematique 2 (1841), 151--159.
\item D.~Hilbert, \emph{Ein Beitrag zur Theorie des Legendre'schen Polynoms}, Acta Mathematica 18 (1894), 155--159;
M.-D.~Choi, \emph{Tricks or treats with the Hilbert matrix}, American Mathematical Monthly 90 (1983), 301--312.
\item P.~Borwein and T.~Erd\'elyi, \emph{The integer Chebyshev problem}, Mathematics of Computation 65 (1996), 661--681.
\item G.~Szeg\H{o}, \emph{Orthogonal Polynomials}, AMS Colloquium Publications 23, 1939 (4th ed.\ 1975), Chapters 2--3.
\item G.~H.~Golub and J.~H.~Welsch, \emph{Calculation of Gauss quadrature rules}, Mathematics of Computation 23 (1969),
221--230: Proposition~\ref{linear:prop:jacobi} as an algorithm.
\item P.~Jordan and J.~von Neumann, \emph{On inner products in linear, metric spaces}, Annals of Mathematics 36 (1935),
719--723.
\item C.~H.~M\"untz, \emph{\"Uber den Approximationssatz von Weierstrass}, Schwarz-Festschrift, Berlin, 1914, 303--312.
\item L.~N.~Trefethen and D.~Bau, \emph{Numerical Linear Algebra}, SIAM, 1997, Lectures 7--11 (QR, Householder, least
squares) and Lecture~19 (stability of least squares).
\item S.~Axler, \emph{Linear Algebra Done Right}, 4th ed., Springer, 2024, Chapters 6--7.
\end{itemize}
